% 张三慧《大学物理学》第4版《静电场中的导体》14.1--14.14
% 原始教材V2 PDF74--75页；配套书精确页码待定位。
% 独立推导候选。图示、非理想模型与版权许可仍待独立复核。
\section*{第14章\quad 静电场中的导体：独立解析}

\paragraph{习题 14.1}
取导体外表面微小面元 \(\Delta S\)，
局部面电荷密度为 \(\sigma\)，其自身无限大薄面近似
对面两侧的场分别为
\(\pm\sigma/(2\varepsilon_0)\boldsymbol n\)，
其中 \(\boldsymbol n\) 是导体外法线。
其他电荷源在面元两侧邻域产生连续的场
\(\boldsymbol E_{\rm other}\)。
导体内部静电平衡给出
\[
 \boldsymbol 0=\boldsymbol E_{\rm other}
 -\frac{\sigma}{2\varepsilon_0}\boldsymbol n
 \quad\Longrightarrow\quad
 \boldsymbol E_{\rm other}
 =\frac{\sigma}{2\varepsilon_0}\boldsymbol n.
\]
外侧两部分相加：
\[
 \boxed{\boldsymbol E_{\rm outside}
 =\frac{\sigma}{\varepsilon_0}\boldsymbol n}.
\]
这是局部边界极限，适用于表面曲率半径远大于
选取面元尺寸的邻域；符号随 \(\sigma\) 取正负。

\paragraph{习题 14.2}
内球半径 \(R_1\)，导体壳内外半径 \(R_2,R_3\)。
设内球未知电荷 \(Q_1\)，外壳已知净电荷 \(Q_2\)。
高斯定律与导体内部场为零，
使壳内表面电荷必为 \(-Q_1\)，
外表面电荷为 \(Q_1+Q_2\)。
无穷远电势为零，则
\[
 \varphi_1=k_e\left[
 Q_1\left(\frac1{R_1}-\frac1{R_2}+\frac1{R_3}\right)
 +\frac{Q_2}{R_3}\right],\quad
 k_e=\frac1{4\pi\varepsilon_0}.
\]
所以可由给定内球电势确定
\[
 \boxed{Q_1=
 \frac{\varphi_1/k_e-Q_2/R_3}
 {1/R_1-1/R_2+1/R_3}}.
\]
球对称场强分段为
\[
 \boxed{E_r(r)=
 \begin{cases}
 0,&0<r<R_1,\\
 k_eQ_1/r^2,&R_1<r<R_2,\\
 0,&R_2<r<R_3,\\
 k_e(Q_1+Q_2)/r^2,&r>R_3.
 \end{cases}}
\]
电势连续且由 \(E_r=-d\varphi/dr\) 给出
\[
 \boxed{\varphi(r)=
 \begin{cases}
 \varphi_1,&r\le R_1,\\
 k_e[Q_1/r-Q_1/R_2+(Q_1+Q_2)/R_3],
   &R_1<r<R_2,\\
 k_e(Q_1+Q_2)/R_3,&R_2\le r\le R_3,\\
 k_e(Q_1+Q_2)/r,&r>R_3.
 \end{cases}}
\]
在各带电薄表面应理解电场为两侧极限，电势保持连续。

\paragraph{习题 14.3}
初态球 \(A\) 总电荷 \(Q_A=3\times10^{-8}\,\mathrm C\)，
壳 \(B\) 总电荷 \(Q_B=2\times10^{-8}\,\mathrm C\)。
同心导体屏蔽要求 \(B\) 内表面
\(Q_{B,\mathrm{in}}=-Q_A=-3\times10^{-8}\,\mathrm C\)，
外表面
\(Q_{B,\mathrm{out}}=Q_A+Q_B=5\times10^{-8}\,\mathrm C\)。
取 \(k_e=8.9876\times10^9\,\mathrm{SI}\)，
内球、外壳电势
\[
 \boxed{\varphi_A=k_e
 [Q_A(1/R_1-1/R_2)+(Q_A+Q_B)/R_3]
 \approx5.62\times10^3\,\mathrm V},
\]
\[
 \boxed{\varphi_B=k_e(Q_A+Q_B)/R_3
 \approx4.49\times10^3\,\mathrm V}.
\]
(2) 首先 \(B\) 接地而 \(A\) 保持隔离带 \(Q_A\)，
此时内表面仍为 \(-Q_A\)，
外表面电荷流至地面为零；
断开地线后 \(B\) 的总电荷固定为 \(-Q_A\)。
随后仅 \(A\) 接地，令其剩余电荷为 \(Q'\)，
则壳内表面 \(-Q'\)，壳外表面 \(Q'-Q_A\)。
因 \(\varphi_A'=0\)，
\[
 0=k_e\!\left[
 Q'\left(\frac1{R_1}-\frac1{R_2}+\frac1{R_3}\right)
 -\frac{Q_A}{R_3}\right].
\]
代入半径 \((0.06,0.08,0.10)\,\mathrm m\)，得
\[
 \boxed{Q'=\frac{Q_A/R_3}{1/R_1-1/R_2+1/R_3}
 \approx2.118\times10^{-8}\,\mathrm C}.
\]
故
\[
 \boxed{Q_{B,\mathrm{in}}'\approx-2.118\times10^{-8}\,\mathrm C},
 \quad
 \boxed{Q_{B,\mathrm{out}}'=Q'-Q_A
 \approx-0.8824\times10^{-8}\,\mathrm C},
\]
\[
 \boxed{\varphi_A'=0,\qquad
 \varphi_B'=k_e(Q'-Q_A)/R_3
 \approx-7.93\times10^2\,\mathrm V}.
\]
先后接地顺序至关重要；若同时接地是不同边界问题。

\paragraph{习题 14.4}
半径为 \(R\) 的接地导体球，
外部点电荷 \(q\) 距球心为 \(r>R\)。
利用唯一性定理，可在球心连线处
距球心 \(a=R^2/r\) 的球内放置
\[
 q_{\rm image}=-q\frac Rr
\]
作为外部电势的等效像电荷，
使球面 \(\varphi(R)=0\)。
把紧贴导体表面的高斯面包含的净感生电荷
等效于球内像电荷总量，
所以
\[
 \boxed{Q_{\rm induced}=-q\frac Rr}.
\]
这里的「像电荷」只是求解外部场的数学工具，
不是导体内真正存在一个点电荷。
若球不接地而保持总电荷为零，
结果会与此不同。

\paragraph{习题 14.5}
三块极宽的平行导体板由上到下记A、B、C，
A与C用导线连为一个等电势导体，
B位于其间，总面电荷密度为
\(\sigma=1.3\times10^{-5}\,\mathrm{C/m^2}\)。
A—B间距 \(d_1=0.050\,\mathrm m\)，
B—C间距 \(d_2=0.030\,\mathrm m\)。
设B上、下两面分别带密度
\(\sigma_{B,t},\sigma_{B,b}\)。
夹层电场由接触面电荷决定，
两夹层电势差的绝对值相同：
\(\sigma_{B,t}d_1=\sigma_{B,b}d_2\)，
并有
\(\sigma_{B,t}+\sigma_{B,b}=\sigma\)。
由此
\[
 \boxed{\sigma_{B,t}=\sigma\frac{d_2}{d_1+d_2}
 =4.875\times10^{-6}\,\mathrm{C/m^2}},\quad
 \boxed{\sigma_{B,b}=\sigma\frac{d_1}{d_1+d_2}
 =8.125\times10^{-6}\,\mathrm{C/m^2}}.
\]
各夹层相邻面的电荷等量异号，
故A的下表面带 \(-\sigma_{B,t}\)，
C的上表面带 \(-\sigma_{B,b}\)。
A、C原来净电荷合计为零，而B净带 \(+\sigma\)；
无穷大、上下无偏置的外侧电场由总电荷确定，
A上表面与C下表面各带 \(+\sigma/2\)。
因此
\[
 \boxed{(\sigma_{A,\mathrm{upper}},\sigma_{A,\mathrm{lower}})
 =(\sigma/2,-\sigma d_2/(d_1+d_2))},
\]
\[
 \boxed{(\sigma_{C,\mathrm{upper}},\sigma_{C,\mathrm{lower}})
 =(-\sigma d_1/(d_1+d_2),\sigma/2)}.
\]
外表面各为 \(6.50\times10^{-6}\,\mathrm{C/m^2}\)。
这个结果基于原题「连接外板的总净电荷为零」
与无限平板近似；有限板边缘电荷分布不均匀。

\paragraph{习题 14.6}
两个内部球形空腔各自中心放置
\(q_b,q_c\)，导体整体本身净电荷为零。
静电平衡时，两个空腔壁感生净电荷分别为
\(-q_b\)、\(-q_c\)，
导体外表面净电荷 \(q_b+q_c\)。
由于空腔球形且点电荷恰在空腔中心，
由唯一性定理各腔内电场
（去除中心点电荷自身场）只有球对称的自感生场，
其在腔中心的矢量合力为零：
\[
 \boxed{\boldsymbol F_b=\boldsymbol0,\quad
 \boldsymbol F_c=\boldsymbol0}.
\]
两者对外部 \(q_d\) 和其他空腔均严格屏蔽，
此结果不依赖 \(q_d\) 与球的距离。
球外远处 \(r\gg R_A\) 的 \(q_d\)
感受到导体外表面净电荷的单极远场，
所以近似
\[
 \boxed{\boldsymbol F_d\simeq
 \frac{q_d(q_b+q_c)}{4\pi\varepsilon_0r^2}
 \hat{\boldsymbol r}},
\]
其中 \(\hat{\boldsymbol r}\) 从导体球心指向 \(q_d\)。
仅最后一个公式是远场近似：外部 \(q_d\) 还会
诱发导体表面极化；当 \(q_b+q_c=0\)
时需保留更高阶极化力，不能宣称力严格为零。

\paragraph{习题 14.7}
在导体表面小面元 \(\Delta S\) 上，
其表面电荷 \(dq=\sigma\Delta S\)。
不能让该面元自己的电场对自身施力；
除自身面元外的其他所有电荷在该处的平均场为
\[
 \boldsymbol E_{\rm other}
 =\frac{\boldsymbol E_{\rm out}
 +\boldsymbol E_{\rm in}}2
 =\frac{\sigma}{2\varepsilon_0}\boldsymbol e_n,
\]
因为导体静电平衡
\(\boldsymbol E_{\rm in}=0\)，
外表面场
\(\boldsymbol E_{\rm out}
 =\sigma\boldsymbol e_n/\varepsilon_0\)。
故每单位面积的力
\[
 \boxed{\boldsymbol f=\frac{dq\boldsymbol E_{\rm other}}
 {\Delta S}
 =\frac{\sigma^2}{2\varepsilon_0}\boldsymbol e_n}.
\]
与电荷符号无关，方向均向导体外部，
单位为 \(\mathrm{N/m^2}\)。

\paragraph{习题 14.8}
金属球壳近似为半径 \(R\) 的孤立导体球，
表面电场与电势关系
\(E(R)=\varphi/R\)。
若气体击穿场强不超过
\(E_{\rm bd}=100\,\mathrm{MV/m}\)，
球壳电势要求 \(\varphi=9.15\,\mathrm{MV}\)，则
\[
 \boxed{R\ge\frac{\varphi}{E_{\rm bd}}
 =0.0915\,\mathrm m}.
\]
维持球面电势不变需以速率
\(\dot Q=320\,\mu\mathrm{C/s}\) 补充电荷，
理想最小输入电功率
\[
 \boxed{P_{\min}=\varphi\dot Q
 =(9.15\times10^6)(320\times10^{-6})
 =2.928\times10^3\,\mathrm W}.
\]
传送带宽 \(w=0.485\,\mathrm m\)，
速率 \(u=33.0\,\mathrm{m/s}\)，
假设只有一面有效输送净电荷，
每秒输送面积为 \(wu\)，故
\[
 \boxed{\sigma_{\rm belt}=\frac{\dot Q}{wu}
 \approx2.00\times10^{-5}\,\mathrm{C/m^2}}.
\]
题明确为「绝缘传送带」，
按两面相对无限大的均匀带电绝缘薄面近似，
每一侧的电场大小为
\[
 \boxed{E_{\rm belt}=\frac{|\sigma_{\rm belt}|}{2\varepsilon_0}
 \approx1.13\times10^6\,\mathrm{V/m}}.
\]
配套辅导书PDF物理页224亦使用该 \(1/2\) 因子。
如果实际带面紧邻有导体并且给定静电平衡边界条件，
才应另外考虑 \(E=\sigma/\varepsilon_0\)
的导体表面表达式；
不能把两种边界条件混用。

\paragraph{习题 14.9}
将接地平面取为 \(z=0\)，
点电荷 \(q\) 位于 \((0,0,h)\)。
在上半空间采用镜像法，
将镜像电荷 \(-q\) 放在
\((0,0,-h)\)，使整个 \(z=0\) 平面等势零。
取平面上距投影点的径向距离为 \(\rho\)，
平面正上方场的垂直分量
\[
 E_z(0^+,\rho)=
 -\frac{2qh}{4\pi\varepsilon_0(\rho^2+h^2)^{3/2}}.
\]
接地金属导体的边界
\(\sigma(\rho)=\varepsilon_0E_z(0^+,\rho)\)，故
\[
 \boxed{\sigma(\rho)=
 -\frac{qh}{2\pi(\rho^2+h^2)^{3/2}}}.
\]
对全平面积分
\(\int_0^\infty\sigma(\rho)2\pi\rho\,d\rho=-q\)，
验证感生总电荷为 \(-q\)。

\paragraph{习题 14.10}
平面接地、点电荷 \(q\) 位于高度 \(z>0\)。
镜像电荷为 \(-q\)，与实际电荷距离 \(2z\)。
真实电荷受到的吸引力大小
\[
 F(z)=\frac{q^2}{4\pi\varepsilon_0(2z)^2}
 =\frac{q^2}{16\pi\varepsilon_0z^2}.
\]
要从 \(h\) 缓慢移到无穷远，
外力必须克服吸引作用，做功
\[
 \boxed{W_{\rm ext}=\int_h^\infty F(z)\,dz
 =\frac{q^2}{16\pi\varepsilon_0h}}.
\]
第一位学生误把像电荷当作
独立固定的真实电荷，得到
\(q^2/(8\pi\varepsilon_0h)\)，
比正确结果大一倍。
镜像电荷位置会随真实电荷移动，
金属感生分布也在变化；
物理系统的静电相互作用能
\(U=\tfrac12 q\varphi_{\rm induced}
 =-q^2/(16\pi\varepsilon_0h)\)
正好验证第二位学生的积分。

\paragraph{习题 14.11}
依原图14.12，把接地平面记为 \(z=0\)，
点电荷 \(q>0\) 位于 \(z=h\)。
由镜像电荷可得距轴投影为 \(\rho\)
处的感生面电荷密度
\[
 \sigma(\rho)=-\frac{qh}
 {2\pi(\rho^2+h^2)^{3/2}}.
\]
从电荷 \(q\) 发出、初始沿\emph{水平方向}
的电场线，把靠近点电荷的一小球面分成上下等
立体角，所围初始通量为 \(q/(2\varepsilon_0)\)。
此通量终止于导体平面内半径为
\(\rho_0\) 的圆盘，故所围感生电荷为
\(-q/2\)。
用高斯定律配合直接积分
\[
 Q_{\rm disk}=\int_0^{\rho_0}
 \sigma(\rho)2\pi\rho\,d\rho
 =-q\left(1-\frac h{\sqrt{\rho_0^2+h^2}}\right)
 =-\frac q2.
\]
因此
\[
 \frac h{\sqrt{\rho_0^2+h^2}}=\frac12,\qquad
 \boxed{\rho_0=\sqrt3\,h}.
\]
电场线与导体交点在点电荷正下方的
投影点向外 \(\sqrt3h\) 的圆周上。
此推导假定图14.12是单个电荷对无限接地平面
的轴对称镜像构型。

\paragraph{习题 14.12}
无限长线电荷 \(\lambda=10^{-8}\,\mathrm{C/m}\)
位于无限大接地导体平面上方
\(h=5.0\,\mathrm m\)，
引入平面下方等距的镜像线电荷
\(-\lambda\)。
其正下方地表场
\[
 \boxed{E_z(0^+)=-\frac{\lambda}{\pi\varepsilon_0h}
 \approx-71.9\,\mathrm{N/C}},
\quad
 \boxed{\sigma(0)=\varepsilon_0E_z(0^+)
 =-\frac{\lambda}{\pi h}
 \approx-6.37\times10^{-10}\,\mathrm{C/m^2}}.
\]
真实线电荷每单位长度受镜像线的吸引力
（两线间距 \(2h\)）大小
\[
 \boxed{\frac FL=
 \frac{\lambda^2}{2\pi\varepsilon_0(2h)}
 =\frac{\lambda^2}{4\pi\varepsilon_0h}
 \approx1.80\times10^{-7}\,\mathrm{N/m}},
\]
方向向地面。
若地面不视作理想接地导体，这只是近似模型。

\paragraph{习题 14.13}
要核验「一滴水中的全部电子」数量，
必须指定滴水质量。
设一滴水的质量约为
\(m_d=0.050\,\mathrm g\)（模型假定），
每个水分子 \(\mathrm{H_2O}\) 含10个电子，
摩尔质量 \(M_w=18.0\,\mathrm{g/mol}\)。
电子总数
\[
 N_e=10\frac{m_d}{M_w}N_A
 \approx1.67\times10^{22},
\]
取走这些电子后地球净增正电荷
\[
 Q=N_ee\approx2.68\times10^3\,\mathrm C.
\]
把地球近似为孤立半径
\(R_\oplus=6.37\times10^6\,\mathrm m\) 的导体球，
其电容
\(C_\oplus=4\pi\varepsilon_0R_\oplus
 \approx7.09\times10^{-4}\,\mathrm F\)。
故电势升高
\[
 \boxed{\Delta V=\frac Q{C_\oplus}
 \approx3.77\times10^6\,\mathrm V},
\]
确为「几百万伏」。
该计算不意味实际能将水滴电子全部移走，
并忽略了大气、离子与地球附近导体的电容修正。

\paragraph{习题 14.14}
总电量 \(q\) 均匀分布在导体球
半径 \(R\) 的外表面，
\(\sigma=q/(4\pi R^2)\)。
静电压力向外、大小
\(p=\sigma^2/(2\varepsilon_0)\)。
取上半球，沿向上对称轴方向的合力为
\[
 F_z=\int_{\rm hemisphere}p\cos\theta\,dA
 =pR^2\int_0^{2\pi}d\phi
 \int_0^{\pi/2}\cos\theta\sin\theta\,d\theta
 =p\pi R^2.
\]
所以
\[
 \boxed{F=\frac{\sigma^2\pi R^2}{2\varepsilon_0}
 =\frac{q^2}{32\pi\varepsilon_0R^2}}.
\]
两半球彼此排斥，上半球向上、下半球向下；
\(q\) 的正负不改变力的方向。
